% ===================================================================
%  图表第 5 号 — 房子和它是怎样盖起来的。
%  Traduction 2026 d'apres texto/io, controlee sur texto/fr, texto/en
%  et texto/es. Consigne : voir texto/zh/10-tubiao-01.tex.
%
%  Le tableau 5 est un DIALOGUE, et les deux volumes composent le nom
%  du parleur en petites capitales : on garde \textsc{}, que le rendu
%  laisse tomber en gardant le texte.
%
%  Les renvois du plan sont des LETTRES — (a) le couloir, (b) le
%  salon... — et non des numeros. On les ecrit comme l'ido les ecrit,
%  y compris la ou l'imprimeur a compose « (1) » pour « (l) ».
% ===================================================================

%%K t05-apar-1 apar
\VUsaut{6.0mm}
\VUtitre{15.0pt}{60}{图表第 5 号}
\VUsaut{1.4mm}
\VUtitre{11.4pt}{0}{房子和它是怎样盖起来的。}
\VUsaut{2.2mm}

%%K t05-tit-1 sub
\VUpk{10.2pt}{40}{凑巧的相遇。}
\VUsaut{3.0mm}

%%K t05-01-1 p
\textsc{约翰}。--- 叔叔，我很想知道一所\VUgras{房子}是怎样盖起来的。

%%K t05-01-2 p
\textsc{叔叔} \textsuperscript{(80)}。--- 这不难，孩子；跟我来，我带你去看
贝尔纳先生 \textsuperscript{(79)} 正在盖的那所，我以后要住在那里。

%%K t05-01-3 p
\textsc{约翰}。--- 那这所房子的\VUgras{图样} \textsuperscript{(76)} 是谁画的？

%%K t05-01-4 p
\textsc{叔叔}。--- 是\VUgras{建筑师} \textsuperscript{(75)}；他画了一张很清楚的
图，通过了，\VUgras{包工头} \textsuperscript{(77)} 就动了工。

%%K t05-01-5 p
\textsc{约翰}。--- 包工头是谁呢？

%%K t05-01-6 p
\textsc{叔叔}。--- 是白先生，你朋友雅各的父亲。他跟建筑师鲁先生一起指挥
工人。

%%K t05-01-7 p
瞧！白先生正好拿着他的\VUgras{尺} \textsuperscript{(78)} 来了，鲁先生也拿着
他的图来了。

%%K t05-01-8 p
两位先生早安。近来可好？

%%K t05-01-9 p
\textsc{鲁先生}。--- 很好，谢谢。约翰早安；这孩子长得真快！

%%K t05-01-10 p
\textsc{叔叔}。--- 可不是，俨然一个小大人了。他很认真；看见什么都要人给他
讲。今天他想知道房子是怎样盖起来的。

%%K t05-tit-2 sub
\VUpk{10.2pt}{40}{工人。}
\VUsaut{3.0mm}

%%K t05-01-11 p
\textsc{白先生}。--- 那就跟我来吧，我的小朋友，我这就讲给你听，因为我很喜欢
肯学的孩子。

%%K t05-01-12 p
你看那个手里拿着\VUgras{铁锹} \textsuperscript{(82)} 的工人：他是\VUgras{挖土工}
\textsuperscript{(81)}。他在挖一条\VUgras{沟} \textsuperscript{(46)}，好埋水管。房子
所在的地也是他挖的，好让瓦工砌起四堵\VUgras{墙}。这些墙埋在土里的那部分
叫\VUgras{地基} \textsuperscript{(2)}。地基围起来的地方就成了\VUgras{地窖}
\textsuperscript{(3)}。这地窖有一扇小窗，叫\VUgras{通风口} \textsuperscript{(5)}，
还有一道\VUgras{门} \textsuperscript{(4)} 和一道楼梯。

%%K t05-01-13 p
\VUgras{瓦工} \textsuperscript{(108)} 接着往上砌墙，留出\VUgras{门}
\textsuperscript{(8)} 和\VUgras{窗} \textsuperscript{(17)} 的洞口。

%%K t05-01-14 p
\textsc{约翰}。--- 可是我看不见他呀，这位瓦工！

%%K t05-01-15 p
\textsc{白先生}。--- 房子上的活他已经做完了。他在那边\VUgras{马厩}
\textsuperscript{(54)} 那儿，站在他的\VUgras{脚手架} \textsuperscript{(110)} 上，
砌\VUgras{洗衣房} \textsuperscript{(56)} 的墙。

%%K t05-01-16 p
他有一把泥刀、一个灰槽和一根\VUgras{铅垂线} \textsuperscript{(109)}。不过这些
名字你记不住。

%%K t05-01-17 p
\textsc{约翰}。--- 我记得住，因为我要向叔叔要一把小泥刀和一个灰槽，铅垂线
我自己用一根绳子做。

%%K t05-tit-3 sub
\VUsaut{3.0mm}
\VUpk{10.2pt}{40}{材料。}
\VUsaut{3.0mm}

%%K t05-01-18 p
\textsc{叔叔}。--- 很好。可是你说说，约翰，盖一所房子要什么？

%%K t05-01-19 p
\textsc{约翰}。--- 要\VUgras{琢好的石头} \textsuperscript{(59)} 和\VUgras{灰浆}
\textsuperscript{(65)}。

%%K t05-01-20 p
\textsc{叔叔}。--- 正是。你看那个戴大帽子的人，坐在他的\VUgras{大车}
\textsuperscript{(136)} 上。他是\VUgras{车夫} \textsuperscript{(135)}。他每天往这里
运\VUgras{石块} \textsuperscript{(60)}。他把车卸在院子里，\VUgras{石匠}
\textsuperscript{(83)} 就把石块琢好，用他们的\VUgras{石锯} \textsuperscript{(85)} 锯
开。一个\VUgras{小工} \textsuperscript{(146)} 把石头搬给瓦工，瓦工再把它们砌上。

%%K t05-01-21 p
这时\VUgras{和灰工} \textsuperscript{(87)} 在调灰浆。他要\VUgras{沙}
\textsuperscript{(62)}、\VUgras{石灰} \textsuperscript{(63)} 和\VUgras{水}
\textsuperscript{(64)}。石灰装在\VUgras{口袋} \textsuperscript{(71)} 里放在他旁边；
一个\VUgras{瓦工帮手} \textsuperscript{(111)} 用\VUgras{手推车} \textsuperscript{(112)}
给他送沙，\VUgras{学徒} \textsuperscript{(113)} 在水泵 \textsuperscript{(58)} 那儿。
他正在装一\VUgras{桶} \textsuperscript{(114)} 水。灰浆一会儿就得。\VUgras{挑灰工}
\textsuperscript{(107)} 接过来，顺着梯子挑到脚手架上，那里有个瓦工正把房子这
一面砌完。

%%K t05-01-22 p
那么，做\VUgras{屋顶} \textsuperscript{(31)} 的工人叫什么？

%%K t05-01-23 p
\textsc{约翰}。--- 他是\VUgras{盖瓦工} \textsuperscript{(118)}。

%%K t05-tit-4 sub
\VUsaut{3.0mm}
\VUpk{10.2pt}{40}{房子的屋顶。}
\VUsaut{3.0mm}

%%K t05-01-24 p
\textsc{白先生}。--- 对，不过先得有\VUgras{木匠} \textsuperscript{(115)}。

%%K t05-01-25 p
木匠做\VUgras{屋架} \textsuperscript{(35)}。他把\VUgras{梁} \textsuperscript{(37)} 接起来，用的是
\VUgras{木栓} \textsuperscript{(117)}。上面那些穿宽绒裤的工人就是木匠。
一个扶着一根小梁，另一个正用\VUgras{斧子} \textsuperscript{(116)} 削木栓。靠近
\VUgras{烟囱} \textsuperscript{(41)} 的那个是盖瓦工。他把板条钉在\VUgras{楞木}
(*)上，再把\VUgras{石板} \textsuperscript{(66)} 钉在板条上。有时他也铺瓦。

%%K t05-noto-1 noto
\VUnotes{143.2mm}{%
(*) \VUgras{Balk-o} = 德语 \textit{Balken}，英语 \textit{joist}，法语
\textit{solive}，意大利语 \textit{travicello}，俄语 \textit{balka}，
西班牙语和葡萄牙语 \textit{viga}。这是一个尚未采用的技术词根，这里提出来
以传达法语 \textit{solive} 的意思，它既不是 \textit{trabo}（法语
\textit{poutre}），也不是小梁。它是一根方木，承着地板和天花板。%
}

%%K t05-01-26 p
马厩的屋顶是\VUgras{瓦的} \textsuperscript{(55)}，房子的屋顶是\VUgras{石板的}
\textsuperscript{(32)}，凉廊的屋顶是\VUgras{锌的} \textsuperscript{(24)}。

%%K t05-01-27 p
\textsc{约翰}。--- 锌屋顶也是盖瓦工做的吗？

%%K t05-01-28 p
\textsc{白先生}。--- 不是，那是\VUgras{白铁工} \textsuperscript{(121)} 或者
\VUgras{管子工}，就是那个正把\VUgras{老虎窗} \textsuperscript{(38)} 装进房顶的
人。\VUgras{檐槽} \textsuperscript{(43)} 和\VUgras{落水管} \textsuperscript{(42)} 也是
他装的。他焊锌，也焊白铁。\VUgras{避雷针} \textsuperscript{(40)} 和
\VUgras{风信鸡} \textsuperscript{(39)} 也是他装上去的，就装在\VUgras{山墙}
\textsuperscript{(36)} 上。

%%K t05-01-29 p
\textsc{约翰}。--- 那么房子盖好了？

%%K t05-tit-5 sub
\VUsaut{3.0mm}
\VUpk{10.2pt}{40}{工具。}
\VUsaut{3.0mm}

%%K t05-01-30 p
\textsc{白先生}。--- 还差一点。\VUgras{泥水匠} \textsuperscript{(99)} 用
\VUgras{砖} \textsuperscript{(61)} 砌隔墙，把房子分成一间间屋子。他给
\VUgras{天花板} \textsuperscript{(26)} 和墙抹上\VUgras{灰泥} \textsuperscript{(70)}。
这会儿他正做\VUgras{凉廊} \textsuperscript{(33)} 的天花板。跟瓦工一样，他也有
\VUgras{泥刀} \textsuperscript{(100)} 和\VUgras{灰槽} \textsuperscript{(101)}。

%%K t05-01-31 p
然后细木工做门、窗、\VUgras{木地板} \textsuperscript{(16)} 和\VUgras{百叶窗}
\textsuperscript{(19)}。

%%K t05-01-32 p
这会儿他在院子里他的\VUgras{工作台} \textsuperscript{(90)} 旁做活。他有一把
\VUgras{锯} \textsuperscript{(94)} 用来锯\VUgras{木头} \textsuperscript{(69)}，还有
\VUgras{刨子} \textsuperscript{(91)}、\VUgras{夹钳} \textsuperscript{(92)}、
\VUgras{凿子} \textsuperscript{(94 \textit{bis})}、\VUgras{圆规} \textsuperscript{(95)}
和一把木\VUgras{槌} \textsuperscript{(95 \textit{bis})}。

%%K t05-01-33 p
\textsc{约翰}。--- 哎呀，做细木工比做瓦工还好玩！叔叔，你给我买一张工作台、
一把锯和一个刨子好吗？我要给小黑做个窝。

%%K t05-01-34 p
\textsc{叔叔}。--- 好的，孩子。

%%K t05-01-35 p
\textsc{约翰}。--- 我真高兴！那站在大门口的人是谁？

%%K t05-tit-6 sub
\VUsaut{3.0mm}
\VUpk{10.2pt}{40}{油漆。}
\VUsaut{3.0mm}

%%K t05-01-36 p
\textsc{白先生}。--- 他是\VUgras{油漆工} \textsuperscript{(102)}。他站在梯子上，
拿着一罐\VUgras{漆} \textsuperscript{(103)} 和他的\VUgras{刷子} \textsuperscript{(104)}。
他在给大门的木头上漆，好让它经久些。

%%K t05-01-37 p
屋里糊墙纸也是他做的。

%%K t05-01-38 p
\VUgras{裱糊工} \textsuperscript{(107)} 站在\VUgras{梯子} \textsuperscript{(106)}
上，在\VUgras{阳台} \textsuperscript{(28)} 那儿装一幅\VUgras{窗帘} \textsuperscript{(29)}。

%%K t05-01-39 p
站在大窗子旁边的工人是\VUgras{玻璃工} \textsuperscript{(96)}。他把\VUgras{玻璃} \textsuperscript{(18)} 嵌上，用的是
\VUgras{油灰} \textsuperscript{(73)}。另一个跪在地窖门
旁的工人是\VUgras{锁匠} \textsuperscript{(97)}。他在装一把\VUgras{锁}
\textsuperscript{(6)}。他的\VUgras{锉} \textsuperscript{(98)} 放在旁边。

%%K t05-01-40 p
\textsc{约翰}。--- 那个用\VUgras{锤子} \textsuperscript{(84)} 打铁的人也是锁匠吗？

%%K t05-01-41 p
\textsc{白先生}。--- 说准确些，他是\VUgras{铁匠} \textsuperscript{(125)}。他在打\VUgras{铁活} \textsuperscript{(67)}，是给
\VUgras{电工} \textsuperscript{(124)} 打的，那电工正在\VUgras{车房}
\textsuperscript{(51)} 顶上。他有一个\VUgras{炉子}
\textsuperscript{(126)}、一个\VUgras{砧} \textsuperscript{(127)} 和一把\VUgras{钳子}
\textsuperscript{(128)}。

%%K t05-01-42 p
电工在装电话的\VUgras{铜线} \textsuperscript{(44)}。

%%K t05-tit-7 sub
\VUpk{10.2pt}{40}{园子里的活。}
\VUsaut{3.0mm}

%%K t05-01-43 p
\textsc{叔叔}。--- \VUgras{院子} \textsuperscript{(45)} 尽头，靠近\VUgras{栅栏}
\textsuperscript{(47)} 和\VUgras{大门} \textsuperscript{(48)} 的地方，你看得见
\VUgras{园丁} \textsuperscript{(129)}。他在整那个小\VUgras{园子} \textsuperscript{(49)}。
他已经用\VUgras{铁锹} \textsuperscript{(131)} 把土翻过了，正在撒种，用\VUgras{耙}
\textsuperscript{(130)} 耙平。回头他用他的\VUgras{喷壶} \textsuperscript{(132)} 浇
\VUgras{花畦} \textsuperscript{(150)}，苗就长起来了。

%%K t05-01-44 p
刚给马梳洗喂料的\VUgras{马夫} \textsuperscript{(133)} 正在擦\VUgras{马车}
\textsuperscript{(52)}，用的是海绵、一个\VUgras{手压泵} \textsuperscript{(134)} 和
一桶水。回头他把车推进
车房，用那根沉重的\VUgras{门闩} \textsuperscript{(53)} 把门插上。

%%K t05-01-45 p
好啦，我想你这下满意了；讲得够多了。

%%K t05-01-46 p
\textsc{约翰}。--- 是的，叔叔，不过我还想看看房子里面。

%%K t05-01-47 p
\textsc{叔叔}。--- 那就明天吧。鲁先生会给你讲图样，告诉你每间屋子叫什么。

%%K t05-tit-8 sub
\VUsaut{3.0mm}
\VUpk{10.2pt}{40}{房子的布置。}
\VUsaut{2.4mm}

%%K t05-apar-2 apar
{\centering\textit{（见图样。）}\par}
\VUsaut{2.4mm}

%%K t05-01-48 p
\textsc{建筑师}。--- 来吧，约翰，我带你把房子的各个部分走一遍。

%%K t05-01-49 p
我们先进\VUgras{底层} \textsuperscript{(7)}。这是\VUgras{门铃}按钮
\textsuperscript{(11)}。我们登上三级\VUgras{梯级} \textsuperscript{(14)}，就是
\VUgras{台阶} \textsuperscript{(9)} 的梯级，跨过\VUgras{门槛} \textsuperscript{(10)}，
就是\VUgras{大门} \textsuperscript{(8)} 的门槛，就到了\VUgras{楼梯} \textsuperscript{(13)} 脚下。

%%K t05-01-50 p
\textsc{约翰}。--- 这是走廊。

%%K t05-01-51 p
\textsc{建筑师}。--- 不，这是\VUgras{门厅} \textsuperscript{(12)}。\VUgras{走廊}
\textsuperscript{(a)} 在里头。右边是\VUgras{客厅} \textsuperscript{(b)} 和
\VUgras{饭厅} \textsuperscript{(c)}；饭厅对面是\VUgras{厨房} \textsuperscript{(d)}
和\VUgras{备餐间} \textsuperscript{(e)}；再往前是\VUgras{书房} \textsuperscript{(f)}。
\VUgras{凉廊} \textsuperscript{(23)} 挨着客厅，它有一道\VUgras{玻璃门}
\textsuperscript{(25)}。

%%K t05-01-52 p
现在我们上楼。扶着\VUgras{扶手} \textsuperscript{(15)}，数一数梯级。一共十四级。
这是\VUgras{楼梯平台} \textsuperscript{(m)}：我们到了二\VUgras{层}
\textsuperscript{(27)}。

%%K t05-tit-9 sub
\VUsaut{3.0mm}
\VUpk{10.2pt}{40}{二层。}
\VUsaut{3.0mm}

%%K t05-01-53 p
楼梯对面是\VUgras{厕所} \textsuperscript{(20)} 和\VUgras{梳妆室} \textsuperscript{(j)}。
右边是一间好\VUgras{卧房} \textsuperscript{(g)}，两扇窗子朝着房子的\VUgras{正面}
\textsuperscript{(1)}。左边是\VUgras{浴室} \textsuperscript{(1)} 和\VUgras{客房} \textsuperscript{(i)}，
客房里有一个大\VUgras{凹间} \textsuperscript{(h)}。卧房旁边是
\VUgras{儿童室} \textsuperscript{(k)}。它通着一个宽\VUgras{露台} \textsuperscript{(21)}。
这露台下面还有一间厕所 \textsuperscript{(20)}，靠墙这里是叫人吃饭时敲的
\VUgras{钟} \textsuperscript{(22)}。

%%K t05-01-54 p
\textsc{约翰}。--- 我们上\VUgras{三层}去吧。

%%K t05-01-55 p
\textsc{建筑师}。--- 没有三层。屋顶底下只有\VUgras{阁楼房} \textsuperscript{(33)}
和顶棚 \textsuperscript{(34)}。我们走完了，可以下去了。

%%K t05-01-56 p
\textsc{约翰}。--- 谢谢您，先生，真有意思。我要把看见的都讲给父亲听。
\VUsaut{4.0mm}
\VUfilet{22mm}
